\documentclass{article}
\usepackage{graphicx} % Required for inserting images
\usepackage[margin=0.75in]{geometry}
\usepackage{amsmath, amssymb, amsfonts}
\usepackage{booktabs}
\usepackage{enumitem}
\usepackage{hyperref}

\newcommand{\softplus}{\operatorname{softplus}}
\newcommand{\softmin}{\operatorname{softmin}}
\newcommand{\sigmoid}{\sigma}


\title{\[\dot{H} Derivation\]}
\author{Sugheerth Sreedharan }
\date{October 4, 2025}

\begin{document}

\maketitle

\section{The CBF of $X_t$}

\subsection*{Notation}

Let $X_t = [w_1, \dots, w_T]^\top \in \mathbb{R}^{2T}$, with each waypoint $w_i = (x_i, y_i)$.
Define the semi-axes of the degree-4 hyperellipsoid obstacle $j$ as $a_j$ and $b_j$. The signed pseudo-distance from waypoint $i$ to obstacle $j$ is defined as:

\begin{equation}
    d_{ij} = \sqrt[4]{\left(\frac{p_{xj} - x_i}{a_j}\right)^4 + \left(\frac{p_{yj} - y_i}{b_j}\right)^4 + \varepsilon\,} - 1
\end{equation}

where the obstacle boundary is evaluated at $1$, and $\varepsilon > 0$ is a regularisation constant that keeps the expression differentiable even when $w_i = p_j$.
Note that $d_{ij} > 0$ means waypoint $i$ is safe from obstacle $j$, and $d_{ij} < 0$ means penetration.

For notational convenience, define:
\begin{equation}
    \rho_{ij} \;=\; \sqrt[4]{\left(\frac{p_{xj} - x_i}{a_j}\right)^4 + \left(\frac{p_{yj} - y_i}{b_j}\right)^4 + \varepsilon\,}
\end{equation}
so that $d_{ij} = \rho_{ij} - 1$.

To focus gradient contributions on violated constraints ($d_{ij} < 0$)
while suppressing safe ones ($d_{ij} > 0$), we apply a smooth asymmetric
transformation via the softplus function:

\begin{equation}
    \tilde{d}_{ij}
    = -c \cdot \softplus\!\left(-\frac{d_{ij}}{c}\right) + \log(2)
    = -c \cdot \log\!\left(1 + e^{-d_{ij}/c}\right) + \log(2)
\end{equation}

where $c > 0$ is a sharpness parameter. This transformation satisfies:
\begin{itemize}
    \item $d_{ij} \gg 0 \;\Rightarrow\; \tilde{d}_{ij} \approx \log(2)$
          \quad (safe obstacles contribute nothing)
    \item $d_{ij} \ll 0 \;\Rightarrow\; \tilde{d}_{ij} \approx d_{ij}$
          \quad (violated obstacles pass through unchanged)
    \item $\tilde{d}_{ij} \leq 0$ for all $d_{ij}$, with smooth transition at zero
\end{itemize}

The waypoint-level CBF uses a softmin over the transformed distances
with temperature $k_1$:

\begin{equation}
    h(w_i)
    = \softmin_j^{k_1}(\tilde{d}_{ij})
    = -k_1 \log \frac{1}{N}\sum_{j=1}^{N} e^{-\tilde{d}_{ij}/k_1}
\end{equation}

The trajectory-level CBF uses a softmin over waypoints with temperature $k_2$:

\begin{equation}
    h(X_t)
    = \softmin_i^{k_2}(h(w_i))
    = -k_2 \log \frac{1}{T}\sum_{i=1}^{T} e^{-h(w_i)/k_2}
\end{equation}

\subsection*{Step 1 --- Outer Softmin Weights (over waypoints, temperature $k_2$)}

To ensure numerical stability when evaluating the exponentials, we factor out the minimum waypoint CBF value. Let $h_{\min} = \min_{i'} h(w_{i'})$.

The stabilized trajectory-level CBF is:

\begin{equation}
    h(X_t)
    = h_{\min} - k_2 \log \left( \frac{1}{T}\sum_{i=1}^{T} e^{-(h(w_i) - h_{\min})/k_2} \right)
\end{equation}

Define the stabilized softmax weights over waypoints:

\begin{equation}
    \sigma_i
    = \frac{e^{-(h(w_i) - h_{\min})/k_2}}{\displaystyle\sum_{i'=1}^{T} e^{-(h(w_{i'}) - h_{\min})/k_2}},
    \qquad \sum_{i=1}^{T} \sigma_i = 1
\end{equation}

By the chain rule, the partial derivative of $h(X_t)$ with respect to
waypoint $w_i$ remains algebraically identical, now utilizing the stabilized weights:

\begin{equation}
    \frac{\partial h(X_t)}{\partial w_i}
    = \sigma_i \cdot \nabla_{w_i} h(w_i)
\end{equation}

\subsection*{Step 2 --- Inner Softmin Weights (over obstacles, temperature $k_1$)}

Similarly, we stabilize the waypoint-level CBF against the most deeply penetrated obstacle. Let $\tilde{d}_{\min, i} = \min_{j'} \tilde{d}_{ij'}$.

The stabilized waypoint-level CBF is:

\begin{equation}
    h(w_i)
    = \tilde{d}_{\min, i} - k_1 \log \left( \frac{1}{N}\sum_{j=1}^{N} e^{-(\tilde{d}_{ij} - \tilde{d}_{\min, i})/k_1} \right)
\end{equation}

Define the stabilized softmax weights over obstacles for waypoint $i$:

\begin{equation}
    \alpha_{ij}
    = \frac{e^{-(\tilde{d}_{ij} - \tilde{d}_{\min, i})/k_1}}{\displaystyle\sum_{j'=1}^{N}
      e^{-(\tilde{d}_{ij'} - \tilde{d}_{\min, i})/k_1}},
    \qquad \sum_{j=1}^{N} \alpha_{ij} = 1
\end{equation}

Since $\tilde{d}_{ij} \leq \log(2)$ for all $i,j$, and the shift guarantees the largest exponent is exactly $0$, all exponential terms are bounded in $(0, 1]$, preventing numerical overflow.

Differentiating $h(w_i)$ via the chain rule yields the same weighted sum, as the constant shift $\tilde{d}_{\min, i}$ vanishes during differentiation:

\begin{equation}
    \nabla_{w_i} h(w_i)
    = \sum_{j=1}^{N} \alpha_{ij} \cdot \nabla_{w_i} \tilde{d}_{ij}
\end{equation}

To compute $\nabla_{w_i} \tilde{d}_{ij}$, we apply the chain rule through
the softplus. Since the derivative of $\softplus$ is the sigmoid
$\sigmoid(\cdot) = (1 + e^{-(\cdot)})^{-1}$, and the $+\log(2)$ offset is constant:

\begin{equation}
    \frac{\partial \tilde{d}_{ij}}{\partial d_{ij}}
    = \sigmoid\!\left(-\frac{d_{ij}}{c}\right)
    = \frac{1}{1 + e^{\,d_{ij}/c}}
\end{equation}

Now we find the gradient of the raw distance with respect to the waypoint. Since
$d_{ij} = \rho_{ij} - 1$, we need $\nabla_{w_i} \rho_{ij}$. Applying the chain rule to the 4th root:

\begin{equation}
    \nabla_{w_i}\, \rho_{ij}
    = \frac{1}{4} \left( \left(\frac{p_{xj} - x_i}{a_j}\right)^4 + \left(\frac{p_{yj} - y_i}{b_j}\right)^4 + \varepsilon\, \right)^{-3/4} \cdot \nabla_{w_i} \left[ \left(\frac{p_{xj} - x_i}{a_j}\right)^4 + \left(\frac{p_{yj} - y_i}{b_j}\right)^4 \right]
\end{equation}

Notice the first term simplifies to $\frac{1}{4\rho_{ij}^3}$. Differentiating the inner components with respect to $x_i$ and $y_i$ respectively:

\begin{equation}
    \frac{\partial}{\partial x_i} \left(\frac{p_{xj} - x_i}{a_j}\right)^4
    = 4 \left(\frac{p_{xj} - x_i}{a_j}\right)^3 \left(-\frac{1}{a_j}\right)
    = -\frac{4(p_{xj} - x_i)^3}{a_j^4}
\end{equation}

The factor of $4$ perfectly cancels the $1/4$ from the root derivative. Therefore:

\begin{equation}
    \frac{\partial d_{ij}}{\partial w_i}
    = -\frac{1}{\rho_{ij}^3} \begin{pmatrix} \frac{(p_{xj} - x_i)^3}{a_j^4} \\ \frac{(p_{yj} - y_i)^3}{b_j^4} \end{pmatrix}
\end{equation}

\textbf{Note:} The 4th root introduces a normalisation by $\rho_{ij}^3$. The regulariser $\varepsilon > 0$ ensures $\rho_{ij} > 0$ everywhere, so this gradient is always well-defined, even at the exact center of the obstacle.

Combining via the chain rule:

\begin{equation}
    \nabla_{w_i} \tilde{d}_{ij}
    = \frac{\partial \tilde{d}_{ij}}{\partial d_{ij}} \cdot \frac{\partial d_{ij}}{\partial w_i}
    = -\frac{1}{\rho_{ij}^3}\,\sigmoid\!\left(-\frac{d_{ij}}{c}\right) \begin{pmatrix} \frac{(p_{xj} - x_i)^3}{a_j^4} \\ \frac{(p_{yj} - y_i)^3}{b_j^4} \end{pmatrix}
\end{equation}

Substituting back:

\begin{equation}
\boxed{
    \nabla_{w_i} h(w_i)
    = -\sum_{j=1}^{N} \frac{\alpha_{ij}}{\rho_{ij}^3}\,
      \sigmoid\!\left(-\frac{d_{ij}}{c}\right) \begin{pmatrix} \frac{(p_{xj} - x_i)^3}{a_j^4} \\ \frac{(p_{yj} - y_i)^3}{b_j^4} \end{pmatrix}
}
\end{equation}

The sigmoid factor $\sigmoid(-d_{ij}/c)$ acts as a \textbf{soft gate}:
\begin{itemize}
    \item $d_{ij} \ll 0$: $\sigmoid(-d_{ij}/c) \to 1$
          \quad (violated obstacles pass full gradient)
    \item $d_{ij} \gg 0$: $\sigmoid(-d_{ij}/c) \to 0$
          \quad (safe obstacles contribute no gradient)
    \item $d_{ij} = 0$:  $\sigmoid(0) = 0.5$
          \quad (smooth transition at the safety boundary)
\end{itemize}

The factor $1/\rho_{ij}^3$ \textbf{normalises} the higher-order displacement terms, making the gradient scale appropriately based on the specific shape and size encoded by $a_j$ and $b_j$.

\subsection*{Step 3 --- Full Gradient}

Since the waypoints $\{w_i\}$ are independent components of $X_t$,
the full gradient is assembled blockwise, utilizing the numerically
stabilized weights $\sigma_i$ and $\alpha_{ij}$ derived above:

\begin{equation}
\boxed{
    \nabla_{X_t} h(X_t)
    = \left\{
        -\sigma_i \sum_{j=1}^{N} \frac{\alpha_{ij}}{\rho_{ij}^3}\,
        \sigmoid\!\left(-\frac{d_{ij}}{c}\right)
        \begin{pmatrix} \frac{(p_{xj} - x_i)^3}{a_j^4} \\ \frac{(p_{yj} - y_i)^3}{b_j^4} \end{pmatrix}
      \right\}_{i=1}^{T}
}
\end{equation}

where $\rho_{ij} = \sqrt[4]{\left(\frac{p_{xj} - x_i}{a_j}\right)^4 + \left(\frac{p_{yj} - y_i}{b_j}\right)^4 + \varepsilon\,}$.

This is a vector in $\mathbb{R}^{T \times 2}$, where the block for waypoint $w_i$
is a triply-weighted sum of normalised degree-4 displacement vectors
from $w_i$ to each obstacle center $p_j$.

\subsection*{Summary of Weights}

\begin{center}
\begin{tabular}{lll}
\toprule
\textbf{Symbol} & \textbf{Parameter} & \textbf{Role} \\
\midrule
$\sigma_i$ & $k_2$ & Importance of waypoint $i$ to trajectory-level CBF \\
$\alpha_{ij}$ & $k_1$ & Importance of obstacle $j$ to waypoint-level CBF \\
$\sigmoid(-d_{ij}/c)$ & $c$ & Gate suppressing safe obstacles from gradient \\
$1/\rho_{ij}^3$ & $\varepsilon, a_j, b_j$ & Normalises the hyperellipsoid displacement; $\varepsilon$ prevents division by zero \\
\bottomrule
\end{tabular}
\end{center}

\paragraph{Effect of parameters.}
\begin{itemize}
    \item $k_1 \to 0$: $\alpha_{ij}$ concentrates on the single most-violated obstacle.
    \item $k_2 \to 0$: $\sigma_i$ concentrates on the single most-dangerous waypoint.
    \item $c \to 0$: sigmoid gate becomes a hard threshold ---
          only $d_{ij} < 0$ obstacles contribute any gradient.
    \item $c \to \infty$: sigmoid gate $\to 0.5$ everywhere ---
          recovers the unmodified softmin gradient (all obstacles contribute equally).
    \item $\varepsilon \to 0$: Recovers the exact hyperellipsoid boundary gradient, but becomes singular at the exact center $w_i = p_j$.
    \item $\varepsilon \to \infty$: $1/\rho_{ij}^3 \to 0$, suppressing all gradient contributions.
\end{itemize}
 
% \subsection*{What Changed from the Original}
 
% The modification is in the definition of $d_{ij}$:
 
% \begin{center}
% \begin{tabular}{lll}
% \toprule
%  & \textbf{Definition} & \textbf{Gradient w.r.t.\ $w_i$} \\
% \midrule
% \textbf{Original} & $d_{ij} = \|w_i - p_j\|_2^2 - r^2 + \gamma\delta$
%     & $2(w_i - p_j)$ \\[6pt]
% \textbf{New} & $d_{ij} = \sqrt{\|w_i - p_j\|_2^2 + \varepsilon\,} - r$
%     & $\dfrac{w_i - p_j}{\rho_{ij}}$ \\
% \bottomrule
% \end{tabular}
% \end{center}
 
% \medskip
 
% Unlike the previous version where $d_{ij}$ was quadratic in $w_i$ (yielding
% a linear gradient), the new $d_{ij}$ involves a square root, which introduces
% a \emph{normalisation factor} $1/\rho_{ij}$. Concretely:
 
% \begin{enumerate}
%     \item The factor of $2$ in the original gradient is replaced by $1/\rho_{ij}$.
%     \item Distant waypoints, which have large $\rho_{ij}$, produce
%           \emph{smaller} gradient magnitudes per unit displacement.
%           In the original formulation, distant waypoints produced
%           \emph{larger} gradients (due to the factor of $2\|w_i - p_j\|$).
%     \item The regulariser $\varepsilon > 0$ ensures the gradient is bounded
%           everywhere, including when $w_i$ coincides with $p_j$.
% \end{enumerate}


\section{Generative Sampling with CBF and Diffusion PF-ODE}

\begin{align}
\frac{dx_t}{dt}
&= f(x_t,t) + \frac{g^2(t)}{2}\, s_\theta(x_t,t)
\end{align}

\subsection{Control Affine Equation}

\begin{align}
\dot{x} = F(x,t) + G(x,t)u
\end{align}
where $F(x,t)$ is the drift and $G(x,t)$ is the control term.
The PF-ODE can be represented as 

\begin{align}
\frac{dx_t}{dt}
&= f(x_t,t) + \frac{g^2(t)}{2}, s_\theta(x_t,t) + I\cdot u
\\
&\text{ with }G(x) = I \text{ and } u=\vec{0} \text{ and } F(x)
= f(x,t) + \frac{g^2(t)}{2} s_\theta(x,t)\\
&\text{ with }G(x) = I \text{ and } u=\vec{0} \text{ and } F(x)
= f(x,t) + \frac{g^2(t)}{2\sigma_t} \epsilon_\theta(x,t)
\end{align}
So, for safe operation, QP objective and constraints can be given by 
\begin{align}
    &\text{minimize } \frac{1}{2}u_d^TQu_d + c^Tu_d + \frac{1}{2}\gamma\mu^2
    \\&\text{S.T. }
    \\&-L_fh(x) -Lgh(x)\cdot u -\alpha(h(x)) -\mu h(x)\le 0
\end{align}
where
\begin{align}
Q = I, \qquad c = \vec{0}
\end{align}

\subsection{Optimal Control formulation using closed form solution of QP with CBF constraint:}

\begin{align}
u^*(x) =
\begin{cases}
-\dfrac{\omega(x)}{d(x)} \nabla h(x), & \omega(x) < 0 \\
0, & \omega(x) \ge 0
\end{cases}
\end{align}

\begin{align}
u^*(x)
=
- Q^{-1}
\left(
c - L_g h(x)^T \lambda(x)
\right)
\end{align}

\begin{align}
\lambda(x) =
\begin{cases}
\dfrac{-\omega(x)}{d(x)}, & x \in \Omega \\
0, & x \notin \Omega
\end{cases}
\end{align}

\begin{align}
\omega(x) =
L_f h(x)
-
L_g h(x) Q^{-1} c
+
\alpha(h(x))
\end{align}

\begin{align}
d(x)
=
L_g h(x) Q^{-1} L_g h(x)^T
+
\gamma^{-1} h(x)^2
\end{align}

\subsection{Substituting assumptions:}
\begin{align}
L_g h(x) = \nabla_x h(x) \cdot I
\end{align}

\begin{align}
u^*(x)
=
\lambda(x)\,\nabla_x h(x)
\end{align}

\begin{align}
\lambda(x)
=
\frac{-\omega(x)}
{\|\nabla_x h(x)\|_2^2 + \gamma^{-1}h(x)^2}
\end{align}

\subsection{Substituting VE ODE Drift:}

\begin{align}
L_f h(x)
=
\nabla_x h(x)^T
\left(
\frac{d \log \alpha_t}{dt} x
-
2\sigma_t^2 \frac{d\lambda_t}{dt}\epsilon_\theta(x,t)
\right)
\end{align}

\begin{align}
\omega(x)
=
\nabla_x h(x)^T
\left(
\frac{d \log \alpha_t}{dt} x
-
2\sigma_t^2 \frac{d\lambda_t}{dt}\epsilon_\theta(x,t)
\right)
+
\alpha(h(x))
\end{align}

\subsection{Final Safe Control Formulation}

\begin{align}
u^*(x)
=
-
\frac{
\nabla_x h(x)^T
\left(
\frac{d \log \alpha_t}{dt} x
-
2\sigma_t^2 \frac{d\lambda_t}{dt}\epsilon_\theta(x,t)
\right)
+
\alpha(h(x))
}
{
\|\nabla_x h(x)\|_2^2
+
\gamma^{-1} h(x)^2
}
\;
\nabla_x h(x)
\end{align}

Integrating
\begin{align}
\dot{x}
=
f(x,t)
+
u^*(x)
\end{align}

using an exponential integrator of the form 
\begin{align}
x_t
&=
e^{\int_s^t f(r)\,dr}\,x_s
+
\int_s^t
e^{\int_\tau^t f(r)\,dr}
\left[
\frac{g^2(\tau)}{2\sigma_\tau}\epsilon_\theta(x_\tau,\tau)
+
u^*(x_\tau,\tau)
\right] d\tau
\end{align}
where,
\begin{align}
u^*(x_\tau,\tau)
&=
\begin{cases}
-\dfrac{
\omega(x_\tau,\tau)
}{
\|\nabla h(x_\tau)\|_2^2+\gamma^{-1}h(x_\tau)^2
}\nabla h(x_\tau),
& \omega(x_\tau,\tau)<0,\\[1.2ex]
0, & \omega(x_\tau,\tau)\ge 0,
\end{cases}
\\[1ex]
\omega(x_\tau,\tau)
&=
\nabla h(x_\tau)^\top
\left(
f(\tau)x_\tau+\dfrac{g^2(\tau)}{2\sigma_\tau}\epsilon_\theta(x_\tau,\tau)
\right)
+\alpha(h(x_\tau)).
\end{align}
the sampling equation can be rewritten as 
\begin{align}
x_t
&=
e^{\int_s^t f(r)\,dr}\,x_s
+
\int_s^t
e^{\int_\tau^t f(r)\,dr}
\left[
\frac{g^2(\tau)}{2\sigma_\tau}\epsilon_\theta(x_\tau,\tau)
-
\frac{
\min\!\left(
0,\,
\nabla h(x_\tau)^\top
\left(
f(\tau)x_\tau+\frac{g^2(\tau)}{2\sigma_\tau}\epsilon_\theta(x_\tau,\tau)
\right)
+\alpha(h(x_\tau))
\right)
}{
\|\nabla h(x_\tau)\|_2^2+\gamma^{-1}h(x_\tau)^2
}
\nabla h(x_\tau)
\right]d\tau.
\end{align}

\subsection{Controlled Probability Flow ODE}

The generative ODE with CBF safety control is:
\begin{align}
\dot{x}_t = f(t)x_t + \frac{g^2(t)}{2\sigma_t}\epsilon_\theta(x_t,t) + u^*(x_t,t)
\end{align}

Define the log-SNR $\lambda_t := \log(\alpha_t/\sigma_t)$. Multiplying through by the integrating factor $e^{-\int f(r)dr} = \alpha_t^{-1}$ and applying variation of constants yields the exact solution:
\begin{align}
x_t = \frac{\alpha_t}{\alpha_s}x_s + \alpha_t\int_s^t \frac{1}{\alpha_\tau}\left[\frac{g^2(\tau)}{2\sigma_\tau}\epsilon_\theta(x_\tau,\tau)\right]d\tau + \alpha_t\int_s^t \frac{1}{\alpha_\tau}\left[ u^*(x_\tau,\tau)\right]d\tau
\end{align}

Changing variable $\tau \to \lambda$ via $d\tau = d\lambda/\dot{\lambda}_\tau$ and using the standard DPM-Solver identity $\frac{g^2(\tau)}{2\sigma_\tau\alpha_\tau\dot{\lambda}_\tau} = -e^{-\lambda}$ yields the exact solution in $\lambda$-space for the score term, keeping the time-integral for the control term:
\begin{align}
x_t = \frac{\alpha_t}{\alpha_s}x_s - \alpha_t\int_{\lambda_s}^{\lambda_t} e^{-\lambda}\epsilon_\theta(\hat{x}_\lambda,\lambda)d\lambda + \alpha_t\int_{s}^{t} \frac{u^*(x_\tau,\tau)}{\alpha_\tau}d\tau
\end{align}

\subsection{Expansion of the Optimal Control $u^*$}

For $g(x) = I$, $Q = I$, $c = 0$, the optimal control to satisfy the CBF is given by:
\begin{align}
u^*(x,\tau) = -\frac{\min(0,\omega(x,\tau))}{\|\nabla h(x)\|_2^2 + \gamma^{-1}h(x)^2}\nabla h(x)
\end{align}
where the boundary condition $\omega$ is defined as:
\begin{align}
\omega(x,\tau) = \nabla h(x)^\top\!\left(f(\tau)x + \frac{g^2(\tau)}{2\sigma_\tau}\epsilon_\theta(x,\tau)\right) + \alpha(h(x))
\end{align}

To evaluate this within a first-order solver, we evaluate the constraint activation strictly at the start of the timestep $s$. Let $\omega_s = \omega(x_s, s)$ and $D_s = \|\nabla h_s\|_2^2 + \gamma^{-1}h_s^2$. 

If $\omega_s < 0$, the constraint is active. Expanding the $u^*$ term by substituting $\omega$ yields:
\begin{align}
u^*(x_\tau,\tau) = -\frac{\nabla h_s}{D_s}\left[ \nabla h_s^\top f(\tau)x_\tau + \nabla h_s^\top \frac{g^2(\tau)}{2\sigma_\tau}\epsilon_\theta(x_\tau,\tau) + \alpha(h_s) \right]
\end{align}

\subsection{Integrate 1st Part as DPM-Solver}

Denote $h_i := \lambda_{t_i} - \lambda_{t_{i-1}} > 0$. For a first-order step (DPM-Solver-1), we approximate the state $x$ and the network prediction $\epsilon_\theta$ as constant over the interval $[\lambda_s, \lambda_t]$ (evaluated at $s = t_{i-1}$).

\paragraph{Score integral (analytic):}
\begin{align}
-\alpha_t\int_{\lambda_s}^{\lambda_t} e^{-\lambda}\epsilon_\theta(x_s,s)d\lambda = -\alpha_t\epsilon_\theta(x_s,s)\left[-e^{-\lambda}\right]_{\lambda_s}^{\lambda_t} = -\sigma_t(e^{h_i}-1)\epsilon_\theta(x_s,s)
\end{align}

\subsection{Integrate 2nd Part (Expanded CBF Control)}

The safe control part of the integral is intractable because the the control depends on the starting state, which itself is evolving over time. Secondly, the control is non-linear because of the min function, which switches between 0 and a valuse based on the $\omega$-function evaluation with state that is evolving over time. 

Therefore, we use the same Taylor Series approximation to freeze the state at time s and split the function into two:

\begin{align}
    \alpha_t\int_{s}^{t} \frac{u^*(x_s,\tau)}{\alpha_\tau}d\tau = 
    \begin{cases}
    0 & \text{if }\omega \ge 0\\
    -\frac{\nabla h_s}{D_s} \left[ (\nabla h_s^\top x_s) \alpha_t\int_s^t \frac{f(\tau)}{\alpha_\tau}d\tau + (\nabla h_s^\top \epsilon_{\theta,s}) \alpha_t\int_s^t \frac{g^2(\tau)}{2\sigma_\tau\alpha_\tau}d\tau + \alpha(h_s) \alpha_t\int_s^t \frac{1}{\alpha_\tau}d\tau \right] & \text{if } \omega < 0
    \end{cases}
\end{align}

 We then integrate the second case. Because we assume $x_\tau \approx x_s$ over the microscopic timestep, the spatial terms ($\nabla h_s$, $x_s$, $\epsilon_{\theta,s}$, $\alpha(h_s)$) are treated as constants and pulled out of the time integrals:

\begin{align}
\alpha_t\int_{s}^{t} \frac{u^*(x_s,\tau)}{\alpha_\tau}d\tau = -\frac{\nabla h_s}{D_s} \left[ (\nabla h_s^\top x_s) \alpha_t\int_s^t \frac{f(\tau)}{\alpha_\tau}d\tau + (\nabla h_s^\top \epsilon_{\theta,s}) \alpha_t\int_s^t \frac{g^2(\tau)}{2\sigma_\tau\alpha_\tau}d\tau + \alpha(h_s) \alpha_t\int_s^t \frac{1}{\alpha_\tau}d\tau \right]
\end{align}

We analytically evaluate these three schedule integrals:
\begin{enumerate}
    \item \textbf{Drift Integral:} Using $f(\tau) = \frac{\dot{\alpha}_\tau}{\alpha_\tau}$, we get $\alpha_t\int_s^t \frac{\dot{\alpha}_\tau}{\alpha_\tau^2}d\tau = \alpha_t\left[-\frac{1}{\alpha_\tau}\right]_s^t = \left(\frac{\alpha_t}{\alpha_s} - 1\right)$
    \item \textbf{Diffusion Integral:} Using the same identity as the score integral, this evaluates to $-\sigma_t(e^{h_i}-1)$
    \item \textbf{Class-K Integral:} Using a constant approximation over the time interval, $\alpha_t\int_s^t \frac{1}{\alpha_\tau}d\tau \approx \frac{\alpha_t}{\alpha_s}(t-s)$
\end{enumerate}

Substituting these back gives the exact-integrated control term for $\omega_s < 0$:
\begin{align}
\alpha_t\int_{s}^{t} \frac{u^*(x_s,\tau)}{\alpha_\tau}d\tau = -\frac{\nabla h_s}{D_s}\left[ \left(\frac{\alpha_t}{\alpha_s} - 1\right)(\nabla h_s^\top x_s) - \sigma_t(e^{h_i}-1)(\nabla h_s^\top \epsilon_{\theta,s}) + \frac{\alpha_t}{\alpha_s}(t-s)\alpha(h_s) \right]
\end{align}

\subsection{Combine the Integrals Together}

Combining the standard score step with the piecewise CBF update ensures we only apply the control when $\omega_s < 0$:

\begin{align}
x_t = 
\begin{cases} 
\frac{\alpha_t}{\alpha_s}x_s - \sigma_t(e^{h_i}-1)\epsilon_{\theta,s} & \text{if } \omega_s \ge 0 \\
\frac{\alpha_t}{\alpha_s}x_s - \sigma_t(e^{h_i}-1)\epsilon_{\theta,s} - \frac{\nabla h_s}{D_s}\left[ \left(\frac{\alpha_t}{\alpha_s} - 1\right)(\nabla h_s^\top x_s) - \sigma_t(e^{h_i}-1)(\nabla h_s^\top \epsilon_{\theta,s}) + \frac{\alpha_t}{\alpha_s}(t-s)\alpha(h_s) \right] & \text{if } \omega_s < 0 
\end{cases}
\end{align}

\subsection{The Final Step (Variance Exploding Assumption)}

For the VE case, the schedule simplifies drastically: $\alpha_t=1$, $f(t) = 0$, and $\sigma_t(e^{h_i}-1) = (\sigma_s - \sigma_t)$. 
Applying these identities, the drift integral $\left(\frac{\alpha_t}{\alpha_s} - 1\right)$ goes to zero, and the scale $\frac{\alpha_t}{\alpha_s}$ becomes 1. 

Evaluating the step from $s = t_{i-1}$ to $t = t_i$, the final exact-integrated update for the active constraint region ($\omega_{t_{i-1}} < 0$) is:

\begin{align}
\boxed{
x_{t_{i}} = x_{t_{i-1}} - (\sigma_{t_{i-1}}-\sigma_{t_i})\epsilon_\theta(x_{t_{i-1}}, t_{i-1}) - \frac{\nabla h_{t_{i-1}}}{\|\nabla h_{t_{i-1}}\|_2^2 + \gamma^{-1}h_{t_{i-1}}^2} \left[ -(\sigma_{t_{i-1}}-\sigma_{t_i})(\nabla h_{t_{i-1}}^\top \epsilon_{\theta}(x_{t_{i-1}}, t_{i-1})) + (t_i-t_{i-1})\alpha(h_{t_{i-1}}) \right]
}
\end{align}
\subsection{Properties}

\begin{itemize}
\item \textbf{Generic:} Valid for any VE/VP schedule — 
      only $\alpha_t$, $\sigma_t$, $\dot\lambda_t$ need to be specified.

\item \textbf{Recovery:} When $\omega_{t_{i-1}} \geq 0$, 
      $\min(0,\omega) = 0$ and the update reduces exactly to 
      standard DPM-Solver-1:
      \begin{align}
          x_{t_{i-1}\to t_i} 
          = \frac{\alpha_{t_i}}{\alpha_{t_{i-1}}}\,x_{t_{i-1}}
          - \sigma_{t_i}(e^{h_i}-1)\,\epsilon_\theta(x_{t_{i-1}},t_{i-1})
      \end{align}

\item \textbf{Consistent order:} Both $\epsilon_\theta$ and $u^*$ 
      are held constant at $x_{t_{i-1}}$, giving uniform 
      $\mathcal{O}(h_i^2)$ truncation error.

\item \textbf{No extra network calls:} The CBF correction 
      reuses $\epsilon_\theta(x_{t_{i-1}}, t_{i-1})$ 
      already computed for the score step.

\item \textbf{Schedule coefficient:}
      $\dfrac{\alpha_{t_i}\,h_i}{\alpha_{t_{i-1}}\,\dot\lambda_{t_{i-1}}}$
      rescales the CBF correction to match the log-SNR timescale, 
      consistently with the $(e^{h_i}-1)$ factor in the score term.
\end{itemize}


\section{AG}

σ(t) is a geometric schedule over discrete levels i = 0 … N:

$$\sigma_i = \sigma_{\min} \left(\frac{\sigma_{\max}}{\sigma_{\min}}\right)^{i/N}$$

Default: σ_min = 0.01, σ_max = 10.0, N = 1000.

The time derivative is:

$$\dot{\sigma}(t) = \sigma(t) \ln!\left(\frac{\sigma_{\max}}{\sigma_{\min}}\right)$$

Training Loss (Denoising Score Matching)
Sample $x_0 \sim p_{\text{data}}$, $\varepsilon \sim \mathcal{N}(0, I)$, pick random level $i$:

$$\mathcal{L} = \mathbb{E}\bigl[|\varepsilon_\theta(x_0 + \sigma_i \varepsilon,\ \sigma_i) - \varepsilon|^2\bigr]$$

The network predicts the noise ε (not the score directly; score = −ε/σ).

DPM-Solver-1 + CBF Update (from image)
One denoising step $t_{i-1} \to t_i$ (σ decreasing):

$$ x_{t_{i-1} \to t_i} = x_{t_{i-1}}
(\sigma_{t_{i-1}} - \sigma_{t_i}),\varepsilon_\theta(x_{t_{i-1}}, t_{i-1})
\frac{h_i ,\sigma_{t_{i-1}}}{\dot\sigma_{t_{i-1}}} \cdot \frac{\min(0,\omega_{t_{i-1}}^{\text{VE}})} {|\nabla h_{t_{i-1}}|2^2 + \gamma^{-1} h{t_{i-1}}^2} ,\nabla h_{t_{i-1}} $$
where

$$ h_i = \log!\left(\frac{\sigma_{t_{i-1}}}{\sigma_{t_i}}\right), \qquad \omega_s^{\text{VE}} = \dot\sigma_s,\nabla h_s^\top \varepsilon_\theta(x_s, s) + \alpha(h_s) $$

$h(x)$ — scalar CBF (positive in safe/free space, negative in walls). We use a differentiable signed-distance-to-walls function built from the maze grid.
$\alpha(h)$ — class-K function; we use the linear $\alpha(h) = \alpha_0 \cdot h$ with $\alpha_0 = 1.0$.
$\gamma$ — regularisation parameter (default 1.0).
The CBF term is zero whenever $\omega^{\text{VE}} \ge 0$ (constraint already satisfied).

\end{document}
